\documentclass[12pt,a4paper]{article}

\usepackage[hidelayout]{../../../Latex/Style/coursework-report-core}
\usepackage{bm}

\addbibresource{res-str.bib}

\title{RES-STR --- Materials, Structural Response and Damage}
\author{}
\date{\today}

\begin{document}

\maketitle

\tableofcontents
\clearpage

\section{Domain Purpose and Scope}
\label{sec:res-str-domain-purpose}

RES-STR records how hydrodynamic pressure, shear, force, impulse and
moment histories from the local fixed-rigid barrier simulations will
be transferred into later structural analyses. The active Studentship
baseline is not a deformable-structure, material-damage, scour or
two-way FSI model. Those topics remain stretch extensions after
hydrodynamic load extraction, data provenance and the V1--V2--V3
validation hierarchy are stable.

\section{Domain Deliverables}
\label{sec:res-str-domain-deliverables}

\input{"RES-STR-MAT - Candidate Materials and Constitutive Properties/res-str-mat.tex"}
% -------------------------------------------------------------------------
% WBS ID: RES-STR-ENV
% Deliverable: Marine Exposure, Degradation and Durability
% Status: Not started
% Evidence status: Not assessed
% Review status: Not reviewed
% Last updated:
% Dependencies:
% Downstream interfaces:
% -------------------------------------------------------------------------

\section{RES-STR-ENV --- Marine Exposure, Degradation and Durability}
\label{sec:res-str-env}

\subsection{Purpose and Scope}
\label{sec:res-str-env-purpose}

% TODO: Define what this Deliverable covers and explicitly excludes.

\subsection{Research Questions}
\label{sec:res-str-env-research-questions}

% TODO: State the questions that must be answered before the Deliverable
% can be accepted.

\subsection{Terminology and Definitions}
\label{sec:res-str-env-definitions}

% TODO: Define the technical terminology, symbols, quantities and
% classifications used in this Deliverable.

\subsection{Literature Evidence}
\label{sec:res-str-env-literature-evidence}

% TODO: Record evidence from peer-reviewed papers, authoritative datasets,
% standards, technical reports and primary documentation.
%
% Do not create a detached literature-summary list. Explain how each source
% supports, contradicts or limits a project decision.

\subsection{Governing Relations and Quantitative Definitions}
\label{sec:res-str-env-governing-relations}

% TODO: Record relevant equations, dimensionless groups, empirical models,
% extraction rules or quantitative definitions.
%
% Use align environments for displayed equations.
% Every displayed equation must have a unique \label{eq:...}.
% Do not place punctuation after displayed equations.

\subsection{Material or Structural System}
\label{sec:res-str-env-material-structural-system}

% TODO: Complete this RES-STR Domain-specific subsection.

\subsection{Load Cases}
\label{sec:res-str-env-load-cases}

% TODO: Complete this RES-STR Domain-specific subsection.

\subsection{Constitutive Behaviour}
\label{sec:res-str-env-constitutive-behaviour}

% TODO: Complete this RES-STR Domain-specific subsection.

\subsection{Response and Stability Measures}
\label{sec:res-str-env-response-stability-measures}

% TODO: Complete this RES-STR Domain-specific subsection.

\subsection{Damage and Failure Modes}
\label{sec:res-str-env-damage-failure-modes}

% TODO: Complete this RES-STR Domain-specific subsection.

\subsection{Fluid--Structure Coupling Requirements}
\label{sec:res-str-env-fluid-structure-coupling-requirements}

% TODO: Complete this RES-STR Domain-specific subsection.

\subsection{Candidate Approaches}
\label{sec:res-str-env-candidate-approaches}

% TODO: Define the credible candidate approaches considered.

\subsection{Critical Comparison}
\label{sec:res-str-env-critical-comparison}

% TODO: Compare candidates using physically and project-relevant criteria.
% Do not select an approach solely on popularity or implementation ease.

\subsection{Assumptions and Applicability}
\label{sec:res-str-env-assumptions}

% TODO: State assumptions and the conditions under which they are valid.

\subsection{Limitations and Failure Conditions}
\label{sec:res-str-env-limitations}

% TODO: State where the evidence, model or selected approach ceases to be
% reliable or applicable.

\subsection{Project-Specific Conclusions}
\label{sec:res-str-env-conclusions}

% TODO: Convert the evidence into conclusions specific to the Studentship
% project. Do not merely restate literature findings.

\subsection{Selected Approach and Rationale}
\label{sec:res-str-env-selected-approach}

% TODO: Record the selected approach only after sufficient evidence exists.
% State the alternatives rejected and the reasons for rejection.

\subsection{Unresolved Decisions}
\label{sec:res-str-env-unresolved-decisions}

% TODO: Record unresolved questions, required evidence and decision owners.

\subsection{Requirements and Handoffs}
\label{sec:res-str-env-requirements}

% TODO: State requirements passed to other Research Domains, Software,
% verification activities or later execution work.

\subsection{Deliverable Completion Record}
\label{sec:res-str-env-completion-record}

% TODO: Record whether the Deliverable is:
%
% - Not started
% - In progress
% - Draft complete
% - Reviewed
% - Accepted
%
% Acceptance requires:
%
% - the research questions to be answered;
% - claims to be supported by appropriate evidence;
% - assumptions and limitations to be explicit;
% - the project-specific conclusion to be stated;
% - downstream requirements to be traceable;
% - unresolved decisions to be closed or formally deferred.

% -------------------------------------------------------------------------
% WBS ID: RES-STR-LOD
% Deliverable: Structural Load Cases and Load Transfer
% Status: In progress
% Evidence status: Hydrodynamic load-transfer requirements established
% Review status: Not reviewed
% Last updated: 2026-07-19
% Dependencies: RES-PHY-LOD, RES-GEO-REP, RES-NUM-SPEC
% Downstream interfaces: RES-STR-FSI, RES-STR-SIM, RES-VER-MET, Software output schema
% -------------------------------------------------------------------------

\section{RES-STR-LOD --- Structural Load Cases and Load Transfer}
\label{sec:res-str-lod}

\subsection{Purpose and Scope}
\label{sec:res-str-lod-purpose}

This Deliverable defines the hydrodynamic load histories that must be
preserved for later structural response, damage or FSI analyses. It
does not activate material damage, deformable barriers, scour or
two-way fluid--structure coupling in the baseline tsunami-barrier
comparison.

\subsection{Research Questions}
\label{sec:res-str-lod-research-questions}

\begin{enumerate}
  \item Which pressure, shear, force, impulse and moment outputs are
    required from the local fluid solver?
  \item How are distributed fluid tractions reduced or mapped without
    losing resultant force and moment?
  \item Which structural quantities are stretch consumers rather than
    active baseline outputs?
\end{enumerate}

\subsection{Terminology and Definitions}
\label{sec:res-str-lod-definitions}

% TODO: Define the technical terminology, symbols, quantities and
% classifications used in this Deliverable.

\subsection{Literature Evidence}
\label{sec:res-str-lod-literature-evidence}

\textcite{baragamage2024fullycoupled} demonstrates the sequence from
three-dimensional tsunami fluid loading to structural motion and
feedback. Source role: supports preserving distributed load histories
for later FSI-capable analysis. Limitation: its reduced bridge
structure does not justify activating full deformable barrier analysis
in the current baseline.

\textcite{istrati2019trappedair} shows that pressure distribution can
alter component demand and moment even where total force is known.
Source role: supports retaining spatial pressure/traction fields, not
only resultant force. Project conclusion: the local hydrodynamic model
shall export complete load histories before any structural extension is
attempted.

\subsection{Governing Relations and Quantitative Definitions}
\label{sec:res-str-lod-governing-relations}

% TODO: Record relevant equations, dimensionless groups, empirical models,
% extraction rules or quantitative definitions.
%
% Use align environments for displayed equations.
% Every displayed equation must have a unique \label{eq:...}.
% Do not place punctuation after displayed equations.

\subsection{Material or Structural System}
\label{sec:res-str-lod-material-structural-system}

% TODO: Complete this RES-STR Domain-specific subsection.

\subsection{Load Cases}
\label{sec:res-str-lod-load-cases}

% TODO: Complete this RES-STR Domain-specific subsection.

\subsection{Hydrodynamic Traction Transfer}
\label{subsec:res_str_lod_hydrodynamic_traction_transfer}

The local three-dimensional fluid model supplies a transient traction field on the wetted structural surface. The traction exerted by the fluid on the solid is

\begin{align}
\bm{t}_{f\rightarrow s}
&=
\bm{\sigma}_{f}\bm{n}_{s}
\label{eq:res_str_lod_fluid_traction}
\end{align}

where \(\bm{n}_{s}\) is the outward unit normal of the current structural surface and the fluid Cauchy stress is

\begin{align}
\bm{\sigma}_{f}
&=
-p\bm{I}+\bm{\tau}_{f}
\label{eq:res_str_lod_fluid_cauchy_stress}
\end{align}

For a finite-element structural discretisation, the distributed interface traction is converted into the structural load vector through

\begin{align}
\bm{f}_{\mathrm{hyd}}(t)
&=
\int_{\Gamma_{fs}(t)}
\bm{N}_{s}^{\mathsf{T}}
\bm{t}_{f\rightarrow s}(t)
\,\mathrm{d}\Gamma
\label{eq:res_str_lod_hydrodynamic_load_vector}
\end{align}

The integration is performed over the current wetted interface \(\Gamma_{fs}(t)\). The transferred field must preserve the resultant force and moment and should preserve interface work when the fluid and structural surface discretisations do not match.

The tsunami FSI model of Baragamage and Wu demonstrates the sequence from three-dimensional fluid loading to structural motion and subsequent feedback to the fluid solver, although its structural representation is reduced rather than a full three-dimensional deformable continuum \autocite{baragamage2024fullycoupled}

\textbf{Selected approach.} Hydrodynamic loading is represented as a
distributed pressure and viscous traction field on the fixed barrier
surface during the active baseline. Reduction to a single resultant
force is not sufficient for later structural or damage studies.
Decision role: \cref{eq:res_str_lod_fluid_traction}--\cref{eq:res_str_lod_hydrodynamic_load_vector}
define the load-transfer contract. Limitation: the equations are a
handoff requirement and do not claim that a structural solver or FSI
run has been executed.

\subsection{Constitutive Behaviour}
\label{sec:res-str-lod-constitutive-behaviour}

% TODO: Complete this RES-STR Domain-specific subsection.

\subsection{Response and Stability Measures}
\label{sec:res-str-lod-response-stability-measures}

% TODO: Complete this RES-STR Domain-specific subsection.

\subsection{Damage and Failure Modes}
\label{sec:res-str-lod-damage-failure-modes}

% TODO: Complete this RES-STR Domain-specific subsection.

\subsection{Fluid--Structure Coupling Requirements}
\label{sec:res-str-lod-fluid-structure-coupling-requirements}

% TODO: Complete this RES-STR Domain-specific subsection.

\subsection{Candidate Approaches}
\label{sec:res-str-lod-candidate-approaches}

% TODO: Define the credible candidate approaches considered.

\subsection{Critical Comparison}
\label{sec:res-str-lod-critical-comparison}

% TODO: Compare candidates using physically and project-relevant criteria.
% Do not select an approach solely on popularity or implementation ease.

\subsection{Assumptions and Applicability}
\label{sec:res-str-lod-assumptions}

% TODO: State assumptions and the conditions under which they are valid.

\subsection{Limitations and Failure Conditions}
\label{sec:res-str-lod-limitations}

% TODO: State where the evidence, model or selected approach ceases to be
% reliable or applicable.

\subsection{Project-Specific Conclusions}
\label{sec:res-str-lod-conclusions}

The Studentship shall first stabilise hydrodynamic load extraction:
pressure, shear, force, moment and impulse histories on fixed barrier
patches. Material deformation and damage remain later consumers of
these histories.

\subsection{Selected Approach and Rationale}
\label{sec:res-str-lod-selected-approach}

The selected near-term approach is one-way hydrodynamic load export
from the local URANS--VOF model to structured output. One-way
CFD-to-FEA and two-way FSI are retained as stretch paths, not baseline
deliverables.

\subsection{Unresolved Decisions}
\label{sec:res-str-lod-unresolved-decisions}

Open decisions are load-mapping tolerances, time-history sampling
rate, structural mesh target, sign convention for moments and the
activation criteria for one-way structural analysis.

\subsection{Requirements and Handoffs}
\label{sec:res-str-lod-requirements}

Software shall emit \texttt{patch\_id}-resolved pressure, shear,
resultant force, impulse, overturning moment, torsional moment,
centre-of-pressure and integration-window metadata. RES-VER-MET shall
define tolerances for force/moment consistency before RES-STR-SIM or
RES-STR-FSI consumes the loads.

\subsection{Deliverable Completion Record}
\label{sec:res-str-lod-completion-record}

% TODO: Record whether the Deliverable is:
%
% - Not started
% - In progress
% - Draft complete
% - Reviewed
% - Accepted
%
% Acceptance requires:
%
% - the research questions to be answered;
% - claims to be supported by appropriate evidence;
% - assumptions and limitations to be explicit;
% - the project-specific conclusion to be stated;
% - downstream requirements to be traceable;
% - unresolved decisions to be closed or formally deferred.

% -------------------------------------------------------------------------
% WBS ID: RES-STR-RSP
% Deliverable: Structural Response and Stability
% Status: Not started
% Evidence status: Not assessed
% Review status: Not reviewed
% Last updated:
% Dependencies:
% Downstream interfaces:
% -------------------------------------------------------------------------

\section{RES-STR-RSP --- Structural Response and Stability}
\label{sec:res-str-rsp}

\subsection{Purpose and Scope}
\label{sec:res-str-rsp-purpose}

% TODO: Define what this Deliverable covers and explicitly excludes.

\subsection{Research Questions}
\label{sec:res-str-rsp-research-questions}

% TODO: State the questions that must be answered before the Deliverable
% can be accepted.

\subsection{Terminology and Definitions}
\label{sec:res-str-rsp-definitions}

% TODO: Define the technical terminology, symbols, quantities and
% classifications used in this Deliverable.

\subsection{Literature Evidence}
\label{sec:res-str-rsp-literature-evidence}

% TODO: Record evidence from peer-reviewed papers, authoritative datasets,
% standards, technical reports and primary documentation.
%
% Do not create a detached literature-summary list. Explain how each source
% supports, contradicts or limits a project decision.

\subsection{Governing Relations and Quantitative Definitions}
\label{sec:res-str-rsp-governing-relations}

% TODO: Record relevant equations, dimensionless groups, empirical models,
% extraction rules or quantitative definitions.
%
% Use align environments for displayed equations.
% Every displayed equation must have a unique \label{eq:...}.
% Do not place punctuation after displayed equations.

\subsection{Three-Dimensional Structural Equation of Motion}
\label{subsec:res_str_rsp_three_dimensional_eom}

For the active hydrodynamic baseline, the barrier is represented in
the fluid model as a fixed rigid solid. The deformable three-dimensional
solid description below is retained as a future structural-response
extension after hydrodynamic load histories have been validated.

Following spatial discretisation, the structural equation of motion is

\begin{align}
\bm{M}_{s}\ddot{\bm{q}}
+
\bm{C}_{s}\dot{\bm{q}}
+
\bm{f}_{\mathrm{int}}\!\left(\bm{q},\bm{z}\right)
&=
\bm{f}_{\mathrm{hyd}}(t)
+
\bm{f}_{\mathrm{body}}(t)
\label{eq:res_str_rsp_structural_eom}
\end{align}

where \(\bm{q}\) contains the structural displacement degrees of freedom, \(\bm{z}\) contains material-history variables, \(\bm{M}_{s}\) is the mass matrix, \(\bm{C}_{s}\) is the damping matrix and \(\bm{f}_{\mathrm{int}}\) is the nonlinear internal-force vector.

The fixed-base boundary condition is

\begin{align}
\bm{d}
&=
\bm{0}
\qquad
\text{on }\Gamma_{\mathrm{base}}
\label{eq:res_str_rsp_fixed_base_boundary}
\end{align}

A nonlinear and irregular tsunami load does not itself require nonlinear material behaviour. Structural nonlinearity arises from finite deformation, nonlinear constitutive response, contact, damage or failure. These mechanisms are deferred until a separate structural activation gate is passed.

\subsection{Updated-Lagrangian Reference Description}
\label{subsec:res_str_rsp_updated_lagrangian}

The Lagrangian formulation governs how the deforming solid is referenced during one tsunami simulation. In an Updated-Lagrangian formulation, the latest converged configuration becomes the reference for the following increment

\begin{align}
\bm{x}_{n+1}
&=
\bm{x}_{n}
+
\Delta\bm{d}_{n+1}
\label{eq:res_str_rsp_updated_lagrangian_position}
\end{align}

This differs from the barrier-class comparison loop. Each future
structural case would begin with its own undeformed design geometry,
while the Updated-Lagrangian formulation governs deformation within
that structural simulation.

If structural response is activated, candidate geometries may undergo
substantial rotation, yielding and permanent shape change. The
Updated-Lagrangian description is therefore retained as the preferred
finite-deformation reference formulation. It is compatible with
three-dimensional transient solid dynamics and incremental
constitutive updates \autocite{nemer2021}

\textbf{Deferred selected approach.} Use an Updated-Lagrangian
three-dimensional finite-deformation solid model only after the
fixed-rigid hydrodynamic load extraction and validation gate is passed.

\subsection{Material or Structural System}
\label{sec:res-str-rsp-material-structural-system}

% TODO: Complete this RES-STR Domain-specific subsection.

\subsection{Load Cases}
\label{sec:res-str-rsp-load-cases}

% TODO: Complete this RES-STR Domain-specific subsection.

\subsection{Constitutive Behaviour}
\label{sec:res-str-rsp-constitutive-behaviour}

% TODO: Complete this RES-STR Domain-specific subsection.

\subsection{Response and Stability Measures}
\label{sec:res-str-rsp-response-stability-measures}

% TODO: Complete this RES-STR Domain-specific subsection.

\subsection{Damage and Failure Modes}
\label{sec:res-str-rsp-damage-failure-modes}

% TODO: Complete this RES-STR Domain-specific subsection.

\subsection{Fluid--Structure Coupling Requirements}
\label{sec:res-str-rsp-fluid-structure-coupling-requirements}

% TODO: Complete this RES-STR Domain-specific subsection.

\subsection{Candidate Approaches}
\label{sec:res-str-rsp-candidate-approaches}

% TODO: Define the credible candidate approaches considered.

\subsection{Critical Comparison}
\label{sec:res-str-rsp-critical-comparison}

% TODO: Compare candidates using physically and project-relevant criteria.
% Do not select an approach solely on popularity or implementation ease.

\subsection{Assumptions and Applicability}
\label{sec:res-str-rsp-assumptions}

% TODO: State assumptions and the conditions under which they are valid.

\subsection{Limitations and Failure Conditions}
\label{sec:res-str-rsp-limitations}

% TODO: State where the evidence, model or selected approach ceases to be
% reliable or applicable.

\subsection{Project-Specific Conclusions}
\label{sec:res-str-rsp-conclusions}

% TODO: Convert the evidence into conclusions specific to the Studentship
% project. Do not merely restate literature findings.

\subsection{Selected Approach and Rationale}
\label{sec:res-str-rsp-selected-approach}

% TODO: Record the selected approach only after sufficient evidence exists.
% State the alternatives rejected and the reasons for rejection.

\subsection{Unresolved Decisions}
\label{sec:res-str-rsp-unresolved-decisions}

% TODO: Record unresolved questions, required evidence and decision owners.

\subsection{Requirements and Handoffs}
\label{sec:res-str-rsp-requirements}

% TODO: State requirements passed to other Research Domains, Software,
% verification activities or later execution work.

\subsection{Deliverable Completion Record}
\label{sec:res-str-rsp-completion-record}

% TODO: Record whether the Deliverable is:
%
% - Not started
% - In progress
% - Draft complete
% - Reviewed
% - Accepted
%
% Acceptance requires:
%
% - the research questions to be answered;
% - claims to be supported by appropriate evidence;
% - assumptions and limitations to be explicit;
% - the project-specific conclusion to be stated;
% - downstream requirements to be traceable;
% - unresolved decisions to be closed or formally deferred.

% -------------------------------------------------------------------------
% WBS ID: RES-STR-DMG
% Deliverable: Damage, Failure and Performance-Loss Models
% Status: In progress
% Evidence status: Damage model deferred; load requirements identified
% Review status: Not reviewed
% Last updated: 2026-07-19
% Dependencies: RES-STR-LOD, RES-STR-RSP, RES-VER-MET
% Downstream interfaces: Future structural damage extension
% -------------------------------------------------------------------------

\section{RES-STR-DMG --- Damage, Failure and Performance-Loss Models}
\label{sec:res-str-dmg}

\subsection{Purpose and Scope}
\label{sec:res-str-dmg-purpose}

This Deliverable records damage and failure modelling as a stretch
extension. The active project baseline is hydrodynamic comparison of
fixed rigid barriers; it shall not infer material damage, scour,
fracture or post-failure debris motion from fluid loads alone.

\subsection{Research Questions}
\label{sec:res-str-dmg-research-questions}

% TODO: State the questions that must be answered before the Deliverable
% can be accepted.

\subsection{Terminology and Definitions}
\label{sec:res-str-dmg-definitions}

% TODO: Define the technical terminology, symbols, quantities and
% classifications used in this Deliverable.

\subsection{Literature Evidence}
\label{sec:res-str-dmg-literature-evidence}

% TODO: Record evidence from peer-reviewed papers, authoritative datasets,
% standards, technical reports and primary documentation.
%
% Do not create a detached literature-summary list. Explain how each source
% supports, contradicts or limits a project decision.

\subsection{Governing Relations and Quantitative Definitions}
\label{sec:res-str-dmg-governing-relations}

% TODO: Record relevant equations, dimensionless groups, empirical models,
% extraction rules or quantitative definitions.
%
% Use align environments for displayed equations.
% Every displayed equation must have a unique \label{eq:...}.
% Do not place punctuation after displayed equations.

\subsection{Material or Structural System}
\label{sec:res-str-dmg-material-structural-system}

% TODO: Complete this RES-STR Domain-specific subsection.

\subsection{Load Cases}
\label{sec:res-str-dmg-load-cases}

% TODO: Complete this RES-STR Domain-specific subsection.

\subsection{Constitutive Behaviour}
\label{sec:res-str-dmg-constitutive-behaviour}

% TODO: Complete this RES-STR Domain-specific subsection.

\subsection{Response and Stability Measures}
\label{sec:res-str-dmg-response-stability-measures}

% TODO: Complete this RES-STR Domain-specific subsection.

\subsection{Damage and Failure Modes}
\label{sec:res-str-dmg-damage-failure-modes}

% TODO: Complete this RES-STR Domain-specific subsection.

\subsection{Initial Failure-Modelling Scope}
\label{subsec:res_str_dmg_initial_failure_scope}

The Updated-Lagrangian formulation permits large deformation but does not independently provide plasticity, damage, fracture or topology change. Earlier structural notes identified this as a possible future route; it is not the active baseline.

Full post-fracture continuation would additionally require crack or damage evolution, contact between separated surfaces, topology change and a fluid treatment for detached structural fragments. These capabilities are outside the active fixed-rigid hydrodynamic baseline.

\textbf{Selected scope.} Defer large nonlinear deformation, material
damage, fragmentation, debris--fluid interaction and scour until the
fixed-rigid hydrodynamic load histories have been validated and
accepted. Source role: \textcite{baragamage2024fullycoupled}
demonstrates that tsunami loading can be coupled to structural motion,
but also supports treating such coupling as a higher-fidelity
extension. Limitation: no project damage model has been selected or
executed.


\subsection{Fluid--Structure Coupling Requirements}
\label{sec:res-str-dmg-fluid-structure-coupling-requirements}

% TODO: Complete this RES-STR Domain-specific subsection.

\subsection{Candidate Approaches}
\label{sec:res-str-dmg-candidate-approaches}

% TODO: Define the credible candidate approaches considered.

\subsection{Critical Comparison}
\label{sec:res-str-dmg-critical-comparison}

% TODO: Compare candidates using physically and project-relevant criteria.
% Do not select an approach solely on popularity or implementation ease.

\subsection{Assumptions and Applicability}
\label{sec:res-str-dmg-assumptions}

% TODO: State assumptions and the conditions under which they are valid.

\subsection{Limitations and Failure Conditions}
\label{sec:res-str-dmg-limitations}

% TODO: State where the evidence, model or selected approach ceases to be
% reliable or applicable.

\subsection{Project-Specific Conclusions}
\label{sec:res-str-dmg-conclusions}

The current Studentship output is a load dataset suitable for later
damage studies, not a damage prediction. Damage claims require a
separate constitutive model, validation evidence and acceptance gate.

\subsection{Selected Approach and Rationale}
\label{sec:res-str-dmg-selected-approach}

The selected approach is formal deferral. Damage metrics may be listed
as future consumers of pressure, shear, force, impulse and moment
histories, but shall not be reported as baseline simulation outputs.

\subsection{Unresolved Decisions}
\label{sec:res-str-dmg-unresolved-decisions}

Open decisions are candidate material families, damage variables,
failure criteria, scour coupling and validation data for damage or
post-failure behaviour.

\subsection{Requirements and Handoffs}
\label{sec:res-str-dmg-requirements}

Software shall not implement damage or deformable-structure logic from
this deliverable. It shall preserve hydrodynamic histories and
geometry/patch identifiers so a later structural extension can consume
them.

\subsection{Deliverable Completion Record}
\label{sec:res-str-dmg-completion-record}

% TODO: Record whether the Deliverable is:
%
% - Not started
% - In progress
% - Draft complete
% - Reviewed
% - Accepted
%
% Acceptance requires:
%
% - the research questions to be answered;
% - claims to be supported by appropriate evidence;
% - assumptions and limitations to be explicit;
% - the project-specific conclusion to be stated;
% - downstream requirements to be traceable;
% - unresolved decisions to be closed or formally deferred.

% -------------------------------------------------------------------------
% WBS ID: RES-STR-FSI
% Deliverable: Fluid–Structure Interaction and Coupling Fidelity
% Status: In progress
% Evidence status: FSI hierarchy retained as stretch extension
% Review status: Not reviewed
% Last updated: 2026-07-19
% Dependencies: RES-STR-LOD, RES-MOD-EQS, RES-NUM-SPEC
% Downstream interfaces: Future FSI extension, RES-VER-MET, Software data schema
% -------------------------------------------------------------------------

\section{RES-STR-FSI --- Fluid–Structure Interaction and Coupling Fidelity}
\label{sec:res-str-fsi}

\subsection{Purpose and Scope}
\label{sec:res-str-fsi-purpose}

This Deliverable records a future FSI capability path and the load
history information needed to support it. The active Studentship
baseline remains fixed terrain and fixed rigid barriers in the local
URANS--VOF model. One-way CFD-to-FEA and two-way FSI are stretch
extensions after hydrodynamic load extraction is stable.

\subsection{Research Questions}
\label{sec:res-str-fsi-research-questions}

% TODO: State the questions that must be answered before the Deliverable
% can be accepted.

\subsection{Terminology and Definitions}
\label{sec:res-str-fsi-definitions}

% TODO: Define the technical terminology, symbols, quantities and
% classifications used in this Deliverable.

\subsection{Literature Evidence}
\label{sec:res-str-fsi-literature-evidence}

% TODO: Record evidence from peer-reviewed papers, authoritative datasets,
% standards, technical reports and primary documentation.
%
% Do not create a detached literature-summary list. Explain how each source
% supports, contradicts or limits a project decision.

\subsubsection{Baragamage and Wu (2024)}
\label{sec:res-str-fsi-evidence-baragamage-2024}

\textcite{baragamage2024fullycoupled} demonstrated a tsunami-impact
framework in which an incompressible free-surface Navier--Stokes model
was coupled with a structural equation of motion. Fluid pressure
generated structural force, while structural displacement and velocity
were returned to the fluid boundary.

The study establishes the feasibility of two-way tsunami FSI but does
not establish that two-way coupling is required for every barrier
candidate. Its bridge representation was structurally reduced, and
parts of the validation used a constrained flow configuration.

\subsubsection{Istrati and Buckle (2019)}
\label{sec:res-str-fsi-evidence-istrati-2019}

\textcite{istrati2019trappedair} demonstrated that the distribution of
pressure and structural demand cannot always be reconstructed from the
total force alone. Changes in internal ventilation altered uplift,
moment, and the distribution of demand between structural components.

\subsection{Governing Relations and Quantitative Definitions}
\label{sec:res-str-fsi-governing-relations}

% TODO: Record relevant equations, dimensionless groups, empirical models,
% extraction rules or quantitative definitions.
%
% Use align environments for displayed equations.
% Every displayed equation must have a unique \label{eq:...}.
% Do not place punctuation after displayed equations.

\subsubsection{Structural State and Momentum Equation}
\label{sec:res-str-fsi-structural-equations}

The structural state is provisionally defined as:

\begin{align}
  \mathcal{Q}_{\mathrm{s}}
  &=
  \left\{
    \mathbf{d}_{\mathrm{s}},
    \dot{\mathbf{d}}_{\mathrm{s}},
    \boldsymbol{\sigma}_{\mathrm{s}},
    \boldsymbol{\xi}_{\mathrm{s}}
  \right\}
  \label{eq:res-str-fsi-structural-state}
\end{align}

where $\boldsymbol{\xi}_{\mathrm{s}}$ contains the internal variables
required by the selected material model.

In continuum form, structural momentum conservation is:

\begin{align}
  \rho_{\mathrm{s}}
  \ddot{\mathbf{d}}_{\mathrm{s}}
  &=
  \nabla_{0}\cdot\mathbf{P}_{\mathrm{s}}
  +
  \rho_{\mathrm{s}}\mathbf{b}_{\mathrm{s}}
  \label{eq:res-str-fsi-structural-momentum}
\end{align}

where $\mathbf{P}_{\mathrm{s}}$ is the first
Piola--Kirchhoff stress tensor.

The corresponding finite-element semi-discrete form is:

\begin{align}
  \mathbf{M}_{\mathrm{s}}
  \ddot{\mathbf{d}}_{\mathrm{s}}
  +
  \mathbf{C}_{\mathrm{s}}
  \dot{\mathbf{d}}_{\mathrm{s}}
  +
  \mathbf{f}_{\mathrm{int}}
  \left(
    \mathbf{d}_{\mathrm{s}},
    \boldsymbol{\xi}_{\mathrm{s}}
  \right)
  &=
  \mathbf{f}_{\mathrm{fluid}}
  +
  \mathbf{f}_{\mathrm{ext}}
  \label{eq:res-str-fsi-semidiscrete-structural-equation}
\end{align}

The formulation shall support more than a single structural
degree of freedom. The initial structural model may be linear elastic,
but the interface shall permit later activation of:

\begin{itemize}
  \item geometric nonlinearity;
  \item material plasticity;
  \item contact and connection behaviour;
  \item damage and failure variables;
  \item rigid-body translation and rotation.
\end{itemize}

\subsection{Fluid--Structure Interaction Model and Activation Strategy}
\label{sec:res-str-fsi-local-coupling}

The structural formulation shall therefore preserve the spatial and
temporal hydrodynamic traction even where the optimisation initially
uses a fixed rigid barrier.

For the current project scope, this statement is a data-preservation
requirement only: fixed-barrier simulations must retain enough surface
history for later structural studies, but they do not move the barrier
or solve structural response.

\subsection{Material or Structural System}
\label{sec:res-str-fsi-material-structural-system}

% TODO: Complete this RES-STR Domain-specific subsection.

\subsection{Load Cases}
\label{sec:res-str-fsi-load-cases}

% TODO: Complete this RES-STR Domain-specific subsection.

\subsection{Constitutive Behaviour}
\label{sec:res-str-fsi-constitutive-behaviour}

% TODO: Complete this RES-STR Domain-specific subsection.

\subsection{Response and Stability Measures}
\label{sec:res-str-fsi-response-stability-measures}

% TODO: Complete this RES-STR Domain-specific subsection.

\subsection{Damage and Failure Modes}
\label{sec:res-str-fsi-damage-failure-modes}

% TODO: Complete this RES-STR Domain-specific subsection.

\subsection{Fluid--Structure Coupling Requirements}
\label{sec:res-str-fsi-fluid-structure-coupling-requirements}

% TODO: Complete this RES-STR Domain-specific subsection.

\subsubsection{Fluid--Structure Interface Conditions}
\label{sec:res-str-fsi-interface-conditions}

At the fluid--structure interface $\Gamma_{\mathrm{fs}}$, kinematic
compatibility requires:

\begin{align}
  \mathbf{u}_{\mathrm{f}}
  &=
  \dot{\mathbf{d}}_{\mathrm{s}}
  \qquad
  \text{on }
  \Gamma_{\mathrm{fs}}
  \label{eq:res-str-fsi-kinematic-compatibility}
\end{align}

Dynamic equilibrium requires:

\begin{align}
  \boldsymbol{\sigma}_{\mathrm{f}}
  \mathbf{n}_{\mathrm{f}}
  +
  \boldsymbol{\sigma}_{\mathrm{s}}
  \mathbf{n}_{\mathrm{s}}
  &=
  \mathbf{0}
  \qquad
  \text{on }
  \Gamma_{\mathrm{fs}}
  \label{eq:res-str-fsi-traction-equilibrium}
\end{align}

The structural load is obtained from the mapped fluid traction:

\begin{align}
  \mathbf{f}_{\mathrm{fluid}}
  &=
  \mathcal{T}_{\mathrm{f}\rightarrow\mathrm{s}}
  \left[
    \boldsymbol{\sigma}_{\mathrm{f}}
    \mathbf{n}_{\mathrm{f}}
  \right]
  \label{eq:res-str-fsi-fluid-to-structure-transfer}
\end{align}

where $\mathcal{T}_{\mathrm{f}\rightarrow\mathrm{s}}$ is the
fluid-to-structure load-transfer operator.

For two-way coupling, the deformed structural state is returned through:

\begin{align}
  \Gamma_{\mathrm{fs}}^{\,n+1}
  &=
  \mathcal{T}_{\mathrm{s}\rightarrow\mathrm{f}}
  \left[
    \mathbf{d}_{\mathrm{s}}^{\,n+1}
  \right]
  \label{eq:res-str-fsi-structure-to-fluid-transfer}
\end{align}

The transfer operators shall conserve resultant force and moment to
within the tolerances defined by `RES-VER-MET`.

\subsubsection{FSI Fidelity Hierarchy}
\label{sec:res-str-fsi-fidelity-hierarchy}

The structural-coupling hierarchy is:

\begin{table}[htbp]
  \centering
  \small
  \renewcommand{\arraystretch}{1.25}
  \begin{tabular}{
      p{0.11\textwidth}
      p{0.25\textwidth}
      p{0.27\textwidth}
      p{0.27\textwidth}
    }
    \hline
    Mode
    &
    Fluid treatment
    &
    Structural treatment
    &
    Intended use
    \\
    \hline
    0
    &
    Local URANS
    &
    Fixed rigid barrier
    &
    Initial fixed-rigid hydrodynamic comparison
    \\
    1
    &
    Local URANS or LES
    &
    One-way CFD-to-FEA
    &
    Stress, deformation, stability, and failure screening
    \\
    2
    &
    Local URANS or LES
    &
    Partitioned two-way FSI
    &
    Cases with material but moderate deformation feedback
    \\
    3
    &
    Local high-fidelity model
    &
    Strongly coupled nonlinear FSI
    &
    Flexible, moving, unstable, or intentionally compliant final
    designs
    \\
    \hline
  \end{tabular}
  \caption{Fluid--structure coupling hierarchy for candidate and
  shortlisted barrier designs}
  \label{tab:res-str-fsi-fidelity-hierarchy}
\end{table}

\subsubsection{Two-Way FSI Activation Gate}
\label{sec:res-str-fsi-activation-gate}

Two-way FSI shall be activated where one or more of the following
conditions are satisfied:

\begin{enumerate}
  \item structural displacement changes the exposed geometry or
    wetted area materially;
  \item structural motion changes a flow passage, opening, or
    overtopping path;
  \item structural velocity is no longer negligible relative to the
    adjacent fluid velocity;
  \item rotation or translation changes the reflected or transmitted
    wave field;
  \item connection flexibility controls the global response;
  \item instability, uplift, sliding, overturning, or large rotation
    is predicted;
  \item the barrier relies on intentional compliant motion for energy
    dissipation;
  \item one-way and two-way pilot simulations produce materially
    different hydrodynamic or structural performance measures.
\end{enumerate}

A preliminary motion-feedback indicator may be defined as:

\begin{align}
  \Pi_{\mathrm{fsi}}
  &=
  \max_{\Gamma_{\mathrm{fs}},t}
  \left(
    \frac{
      \left\lVert
      \dot{\mathbf{d}}_{\mathrm{s}}
      \right\rVert
    }{
      \left\lVert
      \mathbf{u}_{\mathrm{f}}
      \right\rVert
      +
      \epsilon_{\mathrm{u}}
    }
  \right)
  \label{eq:res-str-fsi-motion-feedback-indicator}
\end{align}

where $\epsilon_{\mathrm{u}}$ prevents division by zero.

The threshold for activating two-way FSI shall be determined through
sensitivity testing rather than assigned universally.

\subsection{Reference Descriptions of the Coupled Domains}
\label{subsec:res_str_fsi_reference_descriptions}

The regional and local fluid models use an Eulerian description in which fluid variables are evaluated at spatial positions

\begin{align}
\bm{u}
&=
\bm{u}\!\left(\bm{x},t\right)
\label{eq:res_str_fsi_eulerian_velocity}
\end{align}

The solid uses a Lagrangian material description

\begin{align}
\bm{x}
&=
\bm{\chi}\!\left(\bm{X},t\right)
=
\bm{X}
+
\bm{d}\!\left(\bm{X},t\right)
\label{eq:res_str_fsi_lagrangian_motion}
\end{align}

Impact does not require a Lagrangian fluid description. Retaining an Eulerian fluid mesh is more suitable for breaking, overtopping and free-surface topology change, while the Lagrangian solid description tracks the barrier deformation and material history.

\subsection{Fluid--Solid Interface Conditions}
\label{subsec:res_str_fsi_interface_conditions}

Two-way coupling requires kinematic compatibility and dynamic equilibrium on the moving interface

\begin{align}
\bm{u}_{f}
&=
\dot{\bm{d}}_{s}
\qquad
\text{on }\Gamma_{fs}(t)
\label{eq:res_str_fsi_kinematic_condition}
\\
\bm{\sigma}_{f}\bm{n}_{f}
+
\bm{\sigma}_{s}\bm{n}_{s}
&=
\bm{0}
\qquad
\text{on }\Gamma_{fs}(t)
\label{eq:res_str_fsi_dynamic_condition}
\end{align}

The fluid solver transfers traction to the structure. The structural solver returns the current interface position and velocity to the fluid representation. A one-way calculation omits the return of structural motion, while a two-way calculation repeats this exchange throughout the transient solution.

\subsection{Moving-Boundary Treatment}
\label{subsec:res_str_fsi_moving_boundary_treatment}

A body-fitted Arbitrary Lagrangian--Eulerian method moves the fluid mesh with the structure. The fluid advection velocity is then evaluated relative to the mesh velocity

\begin{align}
\bm{u}_{\mathrm{rel}}
&=
\bm{u}
-
\bm{w}_{m}
\label{eq:res_str_fsi_ale_relative_velocity}
\end{align}

ALE provides a direct conforming interface and conventional near-wall treatment. Its principal limitation for the present application is the risk of fluid-mesh distortion under severe structural deformation, followed by remeshing and transfer of pressure, velocity, VOF and turbulence fields. It remains useful as a reference method for moderate-deformation benchmarks \autocite{girfoglio2021}

If FSI is activated as a stretch extension, the selected reference
family is a sharp immersed structural surface with conservative cut
cells on a fixed Eulerian background mesh. The method avoids global
fluid-mesh deformation while accounting geometrically for the changing
fluid volume and moving solid boundary. Pasquariello et al.
demonstrate conservative finite-volume--finite-element coupling with
cut cells and nonmatching interface transfer for three-dimensional FSI
with large structural deformation \autocite{pasquariello2016}.
Baragamage and Wu apply a related immersed-boundary and cut-cell
strategy to three-dimensional tsunami loading
\autocite{baragamage2024fullycoupled}

\textbf{Deferred selected approach.} Retain a sharp immersed-boundary
representation with conservative cut-cell treatment as the preferred
FSI extension path. Retain ALE as a comparator for
moderate-deformation verification cases.

The selected family must subsequently define cut-cell geometry, changing fluid volumes, small-cell stabilisation, reconstructed wall states, traction recovery and conservative transfer between fluid and structural interfaces.

\subsection{Deferred High-Level FSI Architecture}
\label{subsec:res_str_fsi_current_architecture}

The deferred high-level architecture is

\begin{align}
&\text{Eulerian polyhedral finite-volume two-phase fluid}
\nonumber\\
&\quad+
\text{VOF free-surface capture}
\nonumber\\
&\quad+
\text{Updated-Lagrangian three-dimensional solid FEM}
\nonumber\\
&\quad+
\text{generalised-}\alpha\text{ structural time integration}
\nonumber\\
&\quad+
\text{Newton-type nonlinear structural solution}
\nonumber\\
&\quad+
\text{sharp immersed boundary with conservative cut cells}
\label{eq:res_str_fsi_current_architecture}
\end{align}

The following high-level decisions remain unresolved:

\begin{enumerate}
    \item monolithic or partitioned coupling
    \item loose or strong partitioned coupling if the partitioned family is selected
    \item added-mass stabilisation, relaxation and interface quasi-Newton acceleration
    \item conservative transfer of force, moment, displacement and interface work
    \item coordination of fluid and structural time levels
    \item transition criteria between rigid, one-way and two-way FSI calculations
\end{enumerate}

\subsection{Candidate Approaches}
\label{sec:res-str-fsi-candidate-approaches}

% TODO: Define the credible candidate approaches considered.

\subsection{Critical Comparison}
\label{sec:res-str-fsi-critical-comparison}

% TODO: Compare candidates using physically and project-relevant criteria.
% Do not select an approach solely on popularity or implementation ease.

\subsection{Assumptions and Applicability}
\label{sec:res-str-fsi-assumptions}

% TODO: State assumptions and the conditions under which they are valid.

\subsection{Limitations and Failure Conditions}
\label{sec:res-str-fsi-limitations}

% TODO: State where the evidence, model or selected approach ceases to be
% reliable or applicable.

\subsection{Project-Specific Conclusions}
\label{sec:res-str-fsi-conclusions}

% TODO: Convert the evidence into conclusions specific to the Studentship
% project. Do not merely restate literature findings.

\subsubsection{Project-Specific FSI Decision}
\label{sec:res-str-fsi-project-decision}

The software and data architecture shall be FSI-capable from the
outset. The baseline hydrodynamic comparison model shall nevertheless use a fixed
rigid barrier and retain complete surface traction histories.

The initial structural verification shall use one-way CFD-to-FEA load
transfer. Two-way FSI shall be introduced only where the activation
gate demonstrates that structural motion feeds back materially into
the hydrodynamic response.

This approach retains the equation-level capability for coupled
simulation without imposing fully coupled FSI on the active barrier
comparison campaign.

\subsection{Selected Approach and Rationale}
\label{sec:res-str-fsi-selected-approach}

% TODO: Record the selected approach only after sufficient evidence exists.
% State the alternatives rejected and the reasons for rejection.

\subsection{Unresolved Decisions}
\label{sec:res-str-fsi-unresolved-decisions}

% TODO: Record unresolved questions, required evidence and decision owners.

\subsection{Requirements and Handoffs}
\label{sec:res-str-fsi-requirements}

Software shall preserve load, geometry and patch histories in a form
that a later FSI extension can consume, but shall not implement moving
boundaries, structural time integration or two-way coupling from this
Research update. RES-VER-MET shall define force/moment transfer
consistency metrics before FSI activation.

\subsection{Deliverable Completion Record}
\label{sec:res-str-fsi-completion-record}

% TODO: Record whether the Deliverable is:
%
% - Not started
% - In progress
% - Draft complete
% - Reviewed
% - Accepted
%
% Acceptance requires:
%
% - the research questions to be answered;
% - claims to be supported by appropriate evidence;
% - assumptions and limitations to be explicit;
% - the project-specific conclusion to be stated;
% - downstream requirements to be traceable;
% - unresolved decisions to be closed or formally deferred.

% -------------------------------------------------------------------------
% WBS ID: RES-STR-SIM
% Deliverable: Structural Simulation and Finite-Element Methodology
% Status: Not started
% Evidence status: Not assessed
% Review status: Not reviewed
% Last updated:
% Dependencies:
% Downstream interfaces:
% -------------------------------------------------------------------------

\section{RES-STR-SIM --- Structural Simulation and Finite-Element Methodology}
\label{sec:res-str-sim}

\subsection{Purpose and Scope}
\label{sec:res-str-sim-purpose}

% TODO: Define what this Deliverable covers and explicitly excludes.

\subsection{Research Questions}
\label{sec:res-str-sim-research-questions}

% TODO: State the questions that must be answered before the Deliverable
% can be accepted.

\subsection{Terminology and Definitions}
\label{sec:res-str-sim-definitions}

% TODO: Define the technical terminology, symbols, quantities and
% classifications used in this Deliverable.

\subsection{Literature Evidence}
\label{sec:res-str-sim-literature-evidence}

% TODO: Record evidence from peer-reviewed papers, authoritative datasets,
% standards, technical reports and primary documentation.
%
% Do not create a detached literature-summary list. Explain how each source
% supports, contradicts or limits a project decision.

\subsection{Governing Relations and Quantitative Definitions}
\label{sec:res-str-sim-governing-relations}

% TODO: Record relevant equations, dimensionless groups, empirical models,
% extraction rules or quantitative definitions.
%
% Use align environments for displayed equations.
% Every displayed equation must have a unique \label{eq:...}.
% Do not place punctuation after displayed equations.

\subsection{Material or Structural System}
\label{sec:res-str-sim-material-structural-system}

% TODO: Complete this RES-STR Domain-specific subsection.

\subsection{Load Cases}
\label{sec:res-str-sim-load-cases}

% TODO: Complete this RES-STR Domain-specific subsection.

\subsection{Constitutive Behaviour}
\label{sec:res-str-sim-constitutive-behaviour}

% TODO: Complete this RES-STR Domain-specific subsection.

\subsection{Response and Stability Measures}
\label{sec:res-str-sim-response-stability-measures}

% TODO: Complete this RES-STR Domain-specific subsection.

\subsection{Damage and Failure Modes}
\label{sec:res-str-sim-damage-failure-modes}

% TODO: Complete this RES-STR Domain-specific subsection.

\subsection{Fluid--Structure Coupling Requirements}
\label{sec:res-str-sim-fluid-structure-coupling-requirements}

% TODO: Complete this RES-STR Domain-specific subsection.

\subsection{Structural Spatial and Temporal Discretisation}
\label{subsec:res_str_sim_spatial_temporal_discretisation}

The deferred structural solution stack consists of three-dimensional finite elements, Updated-Lagrangian finite-deformation kinematics, generalised-\(\alpha\) time integration and Newton-type nonlinear equilibrium iterations. The generalised-\(\alpha\) method supplies the temporal discretisation only; it does not replace the finite-element formulation, constitutive law or nonlinear solver. This stack is not part of the active fixed-rigid hydrodynamic baseline.

At each time step, equilibrium is enforced at weighted intermediate states through a residual of the form

\begin{align}
\bm{R}_{s}
&=
\bm{M}_{s}\ddot{\bm{q}}_{n+1-\alpha_{m}}
+
\bm{C}_{s}\dot{\bm{q}}_{n+1-\alpha_{f}}
+
\bm{f}_{\mathrm{int}}\!\left(\bm{q}_{n+1-\alpha_{f}},\bm{z}_{n+1}\right)
-
\bm{f}_{\mathrm{ext},n+1-\alpha_{f}}
\label{eq:res_str_sim_generalised_alpha_residual}
\end{align}

The weighted acceleration and displacement states are

\begin{align}
\ddot{\bm{q}}_{n+1-\alpha_{m}}
&=
\left(1-\alpha_{m}\right)\ddot{\bm{q}}_{n+1}
+
\alpha_{m}\ddot{\bm{q}}_{n}
\label{eq:res_str_sim_weighted_acceleration}
\\
\bm{q}_{n+1-\alpha_{f}}
&=
\left(1-\alpha_{f}\right)\bm{q}_{n+1}
+
\alpha_{f}\bm{q}_{n}
\label{eq:res_str_sim_weighted_displacement}
\end{align}

The displacement and velocity are updated through Newmark-type relations

\begin{align}
\bm{q}_{n+1}
&=
\bm{q}_{n}
+
\Delta t\dot{\bm{q}}_{n}
+
\Delta t^{2}
\left[
\left(\frac{1}{2}-\beta\right)\ddot{\bm{q}}_{n}
+
\beta\ddot{\bm{q}}_{n+1}
\right]
\label{eq:res_str_sim_generalised_alpha_displacement_update}
\\
\dot{\bm{q}}_{n+1}
&=
\dot{\bm{q}}_{n}
+
\Delta t
\left[
\left(1-\gamma\right)\ddot{\bm{q}}_{n}
+
\gamma\ddot{\bm{q}}_{n+1}
\right]
\label{eq:res_str_sim_generalised_alpha_velocity_update}
\end{align}

Under the standard parameter conditions, the method provides second-order accuracy and controllable high-frequency numerical dissipation \autocite{chung1993}. Liu and Marsden demonstrate its use within a three-dimensional nonlinear solid and FSI framework \autocite{liu2018}

\subsection{Nonlinear Equilibrium Solution}
\label{subsec:res_str_sim_nonlinear_equilibrium_solution}

If structural response is activated, finite deformation and nonlinear
constitutive behaviour require an iterative solution at each structural
time step. A Newton correction is obtained from

\begin{align}
\bm{K}_{\mathrm{eff}}^{(k)}
\Delta\bm{q}^{(k)}
&=
-\bm{R}_{s}^{(k)}
\label{eq:res_str_sim_newton_correction}
\\
\bm{q}^{(k+1)}
&=
\bm{q}^{(k)}
+
\Delta\bm{q}^{(k)}
\label{eq:res_str_sim_newton_update}
\end{align}

where \(\bm{K}_{\mathrm{eff}}\) contains the mass, damping and consistent tangent-stiffness contributions. The structural state is accepted only after the residual and displacement correction satisfy the selected nonlinear convergence criteria.

\textbf{Deferred selected approach.} Use generalised-\(\alpha\) time
integration embedded in a Newton-type nonlinear finite-element solution
only after hydrodynamic load extraction is accepted.

\textbf{Unresolved decisions.} The constitutive family, element technology, incompressibility treatment, damping model, spectral-radius parameter and nonlinear convergence tolerances remain to be selected.

\subsection{Candidate Approaches}
\label{sec:res-str-sim-candidate-approaches}

% TODO: Define the credible candidate approaches considered.

\subsection{Critical Comparison}
\label{sec:res-str-sim-critical-comparison}

% TODO: Compare candidates using physically and project-relevant criteria.
% Do not select an approach solely on popularity or implementation ease.

\subsection{Assumptions and Applicability}
\label{sec:res-str-sim-assumptions}

% TODO: State assumptions and the conditions under which they are valid.

\subsection{Limitations and Failure Conditions}
\label{sec:res-str-sim-limitations}

% TODO: State where the evidence, model or selected approach ceases to be
% reliable or applicable.

\subsection{Project-Specific Conclusions}
\label{sec:res-str-sim-conclusions}

% TODO: Convert the evidence into conclusions specific to the Studentship
% project. Do not merely restate literature findings.

\subsection{Selected Approach and Rationale}
\label{sec:res-str-sim-selected-approach}

% TODO: Record the selected approach only after sufficient evidence exists.
% State the alternatives rejected and the reasons for rejection.

\subsection{Unresolved Decisions}
\label{sec:res-str-sim-unresolved-decisions}

% TODO: Record unresolved questions, required evidence and decision owners.

\subsection{Requirements and Handoffs}
\label{sec:res-str-sim-requirements}

% TODO: State requirements passed to other Research Domains, Software,
% verification activities or later execution work.

\subsection{Deliverable Completion Record}
\label{sec:res-str-sim-completion-record}

% TODO: Record whether the Deliverable is:
%
% - Not started
% - In progress
% - Draft complete
% - Reviewed
% - Accepted
%
% Acceptance requires:
%
% - the research questions to be answered;
% - claims to be supported by appropriate evidence;
% - assumptions and limitations to be explicit;
% - the project-specific conclusion to be stated;
% - downstream requirements to be traceable;
% - unresolved decisions to be closed or formally deferred.

% -------------------------------------------------------------------------
% WBS ID: RES-STR-SPEC
% Deliverable: Integrated Structural and Damage Specification
% Status: In progress
% Evidence status: Active structural scope limited to hydrodynamic load handoff
% Review status: Not reviewed
% Last updated: 2026-07-19
% Dependencies: RES-STR-LOD, RES-STR-DMG, RES-STR-FSI
% Downstream interfaces: Future structural extension, Software output schema
% -------------------------------------------------------------------------

\section{RES-STR-SPEC --- Integrated Structural and Damage Specification}
\label{sec:res-str-spec}

\subsection{Purpose and Scope}
\label{sec:res-str-spec-purpose}

This Deliverable consolidates the current structural decision:
baseline simulations treat barriers as fixed rigid geometry and export
hydrodynamic pressure, shear, force, impulse and moment histories.
Material response, deformation, damage, scour and two-way FSI are
stretch extensions requiring separate evidence and validation.

\subsection{Research Questions}
\label{sec:res-str-spec-research-questions}

% TODO: State the questions that must be answered before the Deliverable
% can be accepted.

\subsection{Terminology and Definitions}
\label{sec:res-str-spec-definitions}

% TODO: Define the technical terminology, symbols, quantities and
% classifications used in this Deliverable.

\subsection{Literature Evidence}
\label{sec:res-str-spec-literature-evidence}

% TODO: Record evidence from peer-reviewed papers, authoritative datasets,
% standards, technical reports and primary documentation.
%
% Do not create a detached literature-summary list. Explain how each source
% supports, contradicts or limits a project decision.

\subsection{Governing Relations and Quantitative Definitions}
\label{sec:res-str-spec-governing-relations}

% TODO: Record relevant equations, dimensionless groups, empirical models,
% extraction rules or quantitative definitions.
%
% Use align environments for displayed equations.
% Every displayed equation must have a unique \label{eq:...}.
% Do not place punctuation after displayed equations.

\subsection{Material or Structural System}
\label{sec:res-str-spec-material-structural-system}

% TODO: Complete this RES-STR Domain-specific subsection.

\subsection{Load Cases}
\label{sec:res-str-spec-load-cases}

% TODO: Complete this RES-STR Domain-specific subsection.

\subsection{Constitutive Behaviour}
\label{sec:res-str-spec-constitutive-behaviour}

% TODO: Complete this RES-STR Domain-specific subsection.

\subsection{Response and Stability Measures}
\label{sec:res-str-spec-response-stability-measures}

% TODO: Complete this RES-STR Domain-specific subsection.

\subsection{Damage and Failure Modes}
\label{sec:res-str-spec-damage-failure-modes}

% TODO: Complete this RES-STR Domain-specific subsection.

\subsection{Fluid--Structure Coupling Requirements}
\label{sec:res-str-spec-fluid-structure-coupling-requirements}

% TODO: Complete this RES-STR Domain-specific subsection.

\subsection{Candidate Approaches}
\label{sec:res-str-spec-candidate-approaches}

% TODO: Define the credible candidate approaches considered.

\subsection{Critical Comparison}
\label{sec:res-str-spec-critical-comparison}

% TODO: Compare candidates using physically and project-relevant criteria.
% Do not select an approach solely on popularity or implementation ease.

\subsection{Assumptions and Applicability}
\label{sec:res-str-spec-assumptions}

% TODO: State assumptions and the conditions under which they are valid.

\subsection{Limitations and Failure Conditions}
\label{sec:res-str-spec-limitations}

% TODO: State where the evidence, model or selected approach ceases to be
% reliable or applicable.

\subsection{Project-Specific Conclusions}
\label{sec:res-str-spec-conclusions}

The integrated structural output for the current Studentship is a
load-transfer specification, not a structural-damage model.

\subsection{Selected Approach and Rationale}
\label{sec:res-str-spec-selected-approach}

The selected active approach is fixed-rigid hydrodynamic load
extraction with structured one-way handoff. Deformable structural
simulation and two-way FSI remain deferred.

\subsection{Unresolved Decisions}
\label{sec:res-str-spec-unresolved-decisions}

% TODO: Record unresolved questions, required evidence and decision owners.

\subsection{Requirements and Handoffs}
\label{sec:res-str-spec-requirements}

Software shall preserve load histories, geometry patch identifiers and
integration metadata; it shall not implement damage, structural FEM or
moving-boundary FSI as part of this Research handoff.

\subsection{Deliverable Completion Record}
\label{sec:res-str-spec-completion-record}

% TODO: Record whether the Deliverable is:
%
% - Not started
% - In progress
% - Draft complete
% - Reviewed
% - Accepted
%
% Acceptance requires:
%
% - the research questions to be answered;
% - claims to be supported by appropriate evidence;
% - assumptions and limitations to be explicit;
% - the project-specific conclusion to be stated;
% - downstream requirements to be traceable;
% - unresolved decisions to be closed or formally deferred.


\clearpage
\printbibliography[
  heading=bibintoc,
  title={References}
]

\end{document}
